\par\noindent\textbf{(1)} \(n=2\) 空间态为 \(|200\rangle\) 与 \(|21m\rangle\)（\(m=0,\pm1\)）。算符 \(z=r\cos\theta\) 在宇称下为奇，故同宇称对角元 \(\langle 200|z|200\rangle\)、\(\langle21m|z|21m\rangle\) 为零。角动量选择规则为 \(\Delta l=\pm1,\Delta m=0\)，因此只耦合 \(|200\rangle\) 与 \(|210\rangle\)。
\par\noindent\textbf{(2)} 由标准归一氢原子径向函数直接积分 \( \langle 200|z|210\rangle=-3a_0\)（相位约定可翻转号），其中 \(a_0\) 为Bohr半径。积分可分为角积分 \(\int Y_{00}^*\cos\theta Y_{10}\,d\Omega=1/\sqrt3\) 与径向积分 \(\int R_{20}R_{21}r^3dr=-3\sqrt3a_0\)。
\par\noindent\textbf{(3)} 在 \((|200\rangle,|210\rangle)\) 基底，微扰矩阵为 \(H'_{2\times2}=\begin{pmatrix}0&-3ea_0\mathcal E\\-3ea_0\mathcal E&0\end{pmatrix}\)，本征能修正 \(\Delta E_\pm=\pm3ea_0|\mathcal E|\)，本征态是两态的等幅对称、反对称组合（标签随电场与相位约定）。这是氢样简并能级的线性Stark效应。
\par\noindent\textbf{(4)} \(m=\pm1\) 与 \(2s\) 不耦合，其一阶修正为零；氢基态 \(|100\rangle\) 为非简并且宇称确定，也有 \(\langle100|z|100\rangle=0\)，故其主要Stark响应从二阶开始。真实氢原子精细结构与Lamb位移解除部分简并，过弱场时线性结果的适用性需重新评估。
